\documentclass[10pt]{article}
\usepackage[margin=0.9in]{geometry}
\usepackage{amsmath,amssymb,amsthm,hyperref}
\newtheorem{theorem}{Theorem}
\newcommand{\rise}[2]{(#1)_{#2}}
\newcommand{\step}[2]{(#1)_{#2,2}}
\title{A proof of Maier's Bessel-number formula}
\author{Principia Math in Progress}
\date{13 September 2026}
\begin{document}
\maketitle
\paragraph{Task 1. Prove the generalized Eulerian identity.}
Write \(\rise an=a(a+1)\cdots(a+n-1)\) and
\(\step an=a(a+2)\cdots(a+2n-2)\), with both empty products equal to one.
The generalized Stirling numbers satisfy
\[
 S_{n+1,j}(a,b;r)=(-an+bj+r)S_{n,j}(a,b;r)+S_{n,j-1}(a,b;r),
 \qquad S_{0,0}=1,
\]
and vanish outside \(0\le j\le n\). Put
\(B_{n,j}=S_{n,j}(1,2;0)\) and \(B^{(r)}_{n,j}=S_{n,j}(1,2;r)\).
The generalized Eulerian numbers \(E_{n,k}(a,b;c_0,c_\infty)\)
are similarly normalized and satisfy
\[
 E_{n+1,k}=(-an+bk+c_0)E_{n,k}
             +((a+b)n-b(k-1)+c_\infty)E_{n,k-1}.
 \tag{1}
\]
We prove Conjecture 4.20 of \cite{Maier}, including all its parameters.

\begin{theorem}
Let \(p\ge0\) be an integer, \(\zeta\in\{0,1\}\), and \(m=2p+\zeta\).
For \(0\le k\le n\),
\[
 E_{n,k}(-1,2;c+m,m)=
 \frac{\rise{c+m}n}{\step{c+2p+2\zeta}p}
 \sum_{\ell=0}^p
 \frac{\step{c+1}p}{\step{c+1}{k+\ell}}\,
 B_{m,m-\ell}B^{(m-2\ell)}_{n,n-k}.
 \tag{2}
\]
The equality holds as a rational-function identity in \(c\).
Its right side has removable singularities wherever its displayed
denominators vanish, with polynomial continuation equal to the left side.
\end{theorem}

\begin{proof}
We first work at generic \(c\), so every denominator used below is nonzero.
The Stirling recurrence, or direct coefficient comparison, gives
\[
 \sum_{n\ge j}B^{(r)}_{n,j}\frac{z^n}{n!}
   =\frac{(1+z)^r}{j!}(z+z^2/2)^j.
\]
Consequently
\[
 B_{m,m-\ell}=\frac{m!}{2^\ell\ell!(m-2\ell)!},\qquad
 B^{(r)}_{n,n-k}=\frac{n!}{(n-k)!}
        [z^k](1+z)^r(1+z/2)^{n-k}.                \tag{3}
\]
Let
\[
 P_k(y)=\sum_{\ell=0}^p
   \frac{B_{m,m-\ell}}{\step{c+1}{k+\ell}}\,y^{m-2\ell},
 \quad
 H_{n,k}=\frac{n!}{(n-k)!}[z^k](1+z/2)^{n-k}P_k(1+z),
 \quad
 q=\frac{\step{c+1}p}{\step{c+2p+2\zeta}p}.        \tag{4}
\]
Define \(H_{n,k}=0\) outside \(0\le k\le n\).
The proposed right side of (2) is \(F_{n,k}=q\rise{c+m}n H_{n,k}\).

We record two polynomial identities. Set
\(\alpha=c+2k+m-1\), \(\beta=2m+c+2k-2\), and \(\gamma=m\alpha\).
Then
\[
 \mathcal D_kP_k:=(y^2-1)P_k''-\beta yP_k'+\gamma P_k=0, \tag{5}
\]
and, when \(k\ge1\),
\[
 P_{k-1}=\alpha P_k-yP_k'.                       \tag{6}
\]
To verify (5), write the coefficient of \(y^{m-2\ell}\) in \(P_k\)
as \(a_\ell\). Formula (3) gives
\[
 2(\ell+1)(c+2k+2\ell+1)a_{\ell+1}
    =(m-2\ell)(m-2\ell-1)a_\ell .
\]
These are exactly the coefficient equations in (5);
the top and terminal terms also vanish because \(m-2p\in\{0,1\}\).
For (6), the coefficient multiplier is
\(\alpha-(m-2\ell)=c+2k+2\ell-1\), the ratio of the two
step-two denominators.

We now show
\[
 (n+c+m)H_{n+1,k}
 =(n+2k+c+m)H_{n,k}+(n-2k+2+m)H_{n,k-1}.
 \tag{7}
\]
For \(k=0\), this is immediate from \(H_{n,0}=P_0(1)\).
For \(1\le k\le n+1\), put
\[
 N=n-k+1\ge0,\quad A=n+c+m,\quad d=n-2k+2+m,\quad
 U=1+z/2,\quad y=1+z .
\]
A coefficient calculation using (4) and (6) gives
\begin{align}
 &\frac{N!}{n!}
 \{A H_{n+1,k}-(A+2k)H_{n,k}-dH_{n,k-1}\}\notag\\
 &\hspace{25mm}=-[z^{k-1}]U^N(\mathcal D_kP_k)(1+z)=0.       \tag{8}
\end{align}
Here is an explicit verification of the calculation, to keep (8) self-contained.
For any formal series \(h,P\),
\begin{align*}
 [z^{k-1}]hP'&=k[z^k]hP-[z^{k-1}]h'P,\\
 [z^{k-1}]hP''&=k(k+1)[z^{k+1}]hP
                    -2k[z^k]h'P+[z^{k-1}]h''P.
\end{align*}
Substituting \(P_{k-1}=\alpha P-yP'\), the first expression
in (8) becomes \([z^k]U^{N-1}KP\), where
\[
 K=A(N+k)U-(A+2k)N-d\alpha zU+dkUy-dz(Ny/2+U).
\]
Moving derivatives off \(P\) in \([z^{k-1}]U^N\mathcal D_kP\)
gives \([z^k]U^{N-1}LP\), where
\begin{align*}
 L={}&2k(k-1)U^2-2k(N+1)zU+2(N+1)zU
          +\tfrac12N(N+1)z^2\\
    &-\beta kUy+\beta z(Ny/2+U)+\gamma zU.
\end{align*}
In these expressions \(A=N+k-1+c+m\) and \(d=N-k+1+m\).
Expansion gives \(K+L=0\), proving (8).
This verification also covers \(N=0\), or \(k=n+1\):
the \(H_{n,k}\) term has multiplier \(N=0\), while \(U^{-1}\)
is a well-defined formal series. Thus no undefined factorial is used
at this edge. Equation (5) proves (7).

It remains to check the apex normalization. Using (3),
\begin{align*}
 P_0(1)
 &=\sum_{\ell=0}^p
   \frac{\rise{-p}\ell\,\rise{\frac12-p-\zeta}\ell}
        {\rise{(c+1)/2}\ell\,\ell!}\\
 &=\frac{\rise{c/2+p+\zeta}p}{\rise{(c+1)/2}p}
 =q^{-1}.                                                   \tag{9}
\end{align*}
The middle equality is the finite Chu--Vandermonde identity.
It follows, for example, by expanding \((1-t)^p\) in the beta
integral and then continuing the resulting rational identity in
its parameters. Hence \(F_{0,0}=1\).
Multiplying (7) by \(q\rise{c+m}n\) shows that \(F\) satisfies (1)
with \(a=-1,b=2,c_0=c+m,c_\infty=m\).
The triangular recurrence and its apex determine every entry uniquely,
so \(F=E\) by induction on \(n\).
Finally, this equality at generic \(c\) is an identity of rational
functions. The recurrence makes \(E\) a polynomial in \(c\), proving
the claimed cancellation and continuation at every exceptional value.
\end{proof}

\newpage
\paragraph{Task 2. Find a combinatorial interpretation.}
For parameter values that make all terms nonnegative, explain (2) by
a weighted counting model and interpret the summation index \(\ell\).

\paragraph{Task 3. Study the auxiliary polynomial family.}
Identify \(P_k\) within classical polynomial families, and use its
differential equation to obtain useful connection formulas or zero information.

\paragraph{Task 4. Extend the parameter range of finite formulas.}
Determine which additional pairs \((a,b)\) admit analogous finite
Bessel-type expansions for \(E_{n,k}(a,b;c_0,c_\infty)\).

\paragraph{Task 5. Remove apparent singularities explicitly.}
Derive formulas with manifest polynomial dependence on \(c\), including
efficient evaluations at negative integer and negative odd integer values.

\begin{thebibliography}{1}
\bibitem{Maier}
R.~S.~Maier, \emph{Triangular recurrences, generalized Eulerian numbers,
and related number triangles}, \emph{Advances in Applied Mathematics}
\textbf{146} (2023), 102485.
\href{https://arxiv.org/abs/2207.10224v2}{Accepted manuscript,
arXiv:2207.10224v2}, Definition 4.4, Definition 4.6,
and Conjecture 4.20 (page 44).
\end{thebibliography}
\end{document}
